% !TeX spellcheck = ru_RU
% !TEX root = vkr.tex


\section{Анализ предметной области}
\label{sec:relatedworks}

\subsection{Открытые архитектуры центральных процессоров}

Выбор процессорной архитектуры для учебного стенда, призванного демонстрировать полный цикл разработки и выполнения программ, определяется не только функциональной полнотой, но и степенью доступности спецификаций, инструментария и образовательных ресурсов. Предпочтение отдаётся архитектурам, не обременённым лицензионными ограничениями и допускающим свободное воспроизведение как в симуляции, так и в кремнии. Среди свободно распространяемых ISA выделяются RISC‑V, исторически открытые версии MIPS, OpenRISC и SPARC V8.

Архитектура MIPS на протяжении десятилетий служила академическим целям и до сих пор встречается в учебных курсах по организации ЭВМ \cite{MIPS}. Однако развитие её открытых вариантов фактически остановилось, а современные реализации требуют приобретения коммерческих лицензий, что противоречит принципу воспроизводимости. Семейство ARM, доминирующее во встраиваемых системах, остаётся проприетарным, и его свободное использование ограничено лицензионными соглашениями \cite{ARM}.

Архитектура RISC‑V лишена названных ограничений \cite{RISCV}. Её спецификации распространяются под свободными лицензиями, разрешающими любые реализации без отчислений. Модульная структура ISA позволяет ограничиться базовым подмножеством целочисленных инструкций RV32I, которое сочетает компактность, достаточную для понимания студентами, с полноценной поддержкой современных RISC‑принципов: фиксированная длина инструкции, регистровая архитектура, простые способы адресации. Зрелая экосистема (GCC, LLVM, эмуляторы, отладчики) и активное сообщество делают RISC‑V обоснованным выбором для создания открытого учебного процессора.

\subsection{Легковесные UNIX‑подобные операционные системы}

Операционная система (ОС) представляет собой комплекс программного обеспечения, управляющий аппаратными ресурсами вычислительной установки и предоставляющий интерфейс взаимодействия между пользователем и устройством. К её фундаментальным функциям относятся: загрузка и исполнение прикладных программ с образованием процессов; распределение процессорного времени между вычислительными задачами (планирование); управление оперативной памятью, включая механизмы виртуальной памяти, позволяющие расширять адресное пространство за счёт дисковой подсистемы; организация файловой системы и регламентация доступа к данным на постоянных носителях; взаимодействие с периферийным оборудованием через драйверы; предоставление пользовательского интерфейса — командной строки или графической среды. В учебном стенде, нацеленном на демонстрацию цикла работы программ, операционная система должна не только реализовывать перечисленные функции, но и обладать обозримым, хорошо документированным исходным кодом, допускающим детальный анализ в рамках семестрового курса.

Традиционные UNIX‑подобные системы общего назначения, такие как GNU/Linux, не удовлетворяют последнему требованию. Современное ядро Linux содержит миллионы строк кода, а его изучение затруднено даже на уровне ранних версий. Первая версия ядра (0.01), выпущенная Линусом Торвальдсом в сентябре 1991 г. \cite{fstlinux}, являлась прототипом с базовой поддержкой управления памятью, многозадачности, элементарной файловой системы и взаимодействия через терминал. Однако её документация оставалась фрагментарной, а сообщество разработчиков в тот период было ориентировано скорее на эксперименты с кодом, чем на его систематическое описание. Компактные сборки Linux с использованием BusyBox \cite{BusyBox} не решают проблему: BusyBox объединяет множество стандартных утилит в одном исполняемом файле, скрадывая границы между отдельными командами и затрудняя понимание их взаимодействия.

Ряд исторических и специализированных систем также не пригоден для образовательных целей. Операционная система CP/M \cite{CPM}, созданная в 1973 г. на языке PL/M для процессоров Intel 8080 и Zilog Z80, стала фактическим стандартом 8‑битных микрокомпьютеров, однако её архитектура безнадёжно устарела, а оригинальные материалы труднодоступны. Отечественная MicroDOS \cite{MicroDOS} для персональных компьютеров БК‑0010 не поддерживается с 1992 г., имеет закрытый исходный код и лишена возможности модификации. Embox \cite{Embox} представляет собой ОС реального времени, нацеленную на встроенные системы с ограниченными ресурсами. Глубокая привязка к конкретным аппаратным платформам и специфика задач реального времени не позволяют на её основе изучать универсальные механизмы управления процессами и файловыми подсистемами. Проект GNU Mes \cite{GNUMes} решает задачу самодостаточной сборки компиляторов (бутстрэппинга) и предоставляет интерпретатор Scheme вместе с простым компилятором C, однако не содержит многих компонентов полноценной ОС — подсистем управления процессами, сетевого стека и т. п.

Шестая редакция операционной системы Unix (Unix V6), выпущенная Bell Labs в 1975 г. для архитектуры PDP‑11 \cite{V6}, на протяжении десятилетий служила учебным пособием. Благодаря небольшому объёму (порядка 10 000 строк) и ясной структуре она подходила для изучения принципов проектирования операционных систем. Вместе с тем код написан на раннем диалекте языка C и ориентирован на устаревшее оборудование; отсутствуют поддержка многопоточности и развитые механизмы управления памятью, что ограничивает перенос получаемых знаний на современные платформы.

Указанные противоречия разрешает учебная операционная система Xv6, разработанная в Массачусетском технологическом институте \cite{Xv6}. Xv6 воспроизводит архитектуру Unix V6, однако полностью переписана на ANSI C, снабжена подробными комментариями, поддерживает многопроцессорность и адаптирована для архитектур x86 и RISC‑V. Объём исходного кода составляет около 9000 строк — величина, допускающая детальное изучение в течение нескольких академических занятий. Система реализует защищённое ядро, виртуальную память, файловую систему и стандартный интерфейс системных вызовов, а благодаря запуску в эмуляторах (например, QEMU) становится удобным объектом трассировки и отладки. Некоторым ограничением можно считать отсутствие поддержки исполнения на физическом оборудовании, однако для учебного стенда, ориентированного на Verilog‑модель процессора, эмуляционный характер исполнения является скорее преимуществом.

Альтернативные учебные ОС, такие как Minix с микроядерной архитектурой \cite{Minix}, или uClinux, нацеленная на системы без аппаратной поддержки виртуальной памяти, либо добавляют избыточную концептуальную сложность, либо не охватывают весь спектр механизмов, характерных для полноценного UNIX‑ядра. Таким образом, Xv6, сочетающая простоту, полноту и прямую поддержку RISC‑V, выступает более подходящей операционной системой для проектируемого стенда.

\subsection{Компактные компиляторы языка Си}

Включение в стенд компилятора, обеспечивающего трансляцию прикладных программ в код целевой архитектуры, необходимо для реализации сценария «исходный код — исполняемый образ — симуляция». Компилятор должен, с одной стороны, поддерживать стандарт языка Си, с другой — обладать умеренным объёмом исходного кода, допускающим адаптацию бэкенда. Промышленные пакеты GCC и LLVM этому требованию не удовлетворяют: их внутреннее устройство чрезвычайно сложно, а модификация генератора кода для новой, пусть и похожей, архитектуры трудоёмка.

Среди легковесных компиляторов выделяются два кандидата: LCC (Little C Compiler) и TinyCC. LCC  в течение многих лет применялся в образовании благодаря ясной архитектуре и чёткому разделению стадий компиляции, однако развитие проекта остановилось \cite{LCC}, и генератор кода для RISC‑V в нём отсутствует. TinyCC (Tiny C Compiler) \cite{TinyCC}, напротив, продолжает поддерживаться, реализован на языке Си, имеет сравнительно небольшой размер (около 200 строк) и изначально проектировался с прицелом на высокую скорость компиляции и возможность встраивания. Гибкая система описания целевых платформ позволила относительно легко добавить поддержку архитектуры RISC‑V, что и было выполнено в рамках данной работы. Совокупность этих свойств делает TinyCC соответствующим критериям простоты и расширяемости.

\subsection{Учебные компьютерные платформы: от Nand2Tetris до FPGA‑курсов}

Учебные компьютеры представляют собой виртуальные или синтезируемые модели вычислительных устройств, предназначенные для освоения принципов архитектуры ЭВМ без необходимости приобретения физического оборудования. При анализе таких платформ целесообразно учитывать несколько ключевых характеристик.

Прежде всего, важна организация набора инструкций (ISA): принадлежность к классу RISC, CISC или оригинальной архитектуре; разрядность (8, 16 или 32 бита); набор поддерживаемых операций — арифметических, логических, операций с памятью, переходов; доступные режимы адресации и наличие системного стека. Эти свойства непосредственно определяют сложность модели и степень её приближения к реальным процессорам.

Состав компонентов также влияет на полноту отражения аппаратных механизмов. Типичный учебный компьютер включает программный счётчик, регистр состояния с флагами, стек, арифметико-логическое устройство (ALU) и блок управления. Более развитые модели могут содержать многоступенчатый конвейер, кэш-память, средства моделирования прерываний, таймеров, контроллеров прямого доступа к памяти (DMA) и даже элементы многоядерности. Совокупность этих компонентов определяет, насколько детально воспроизводятся конвейеризация, иерархия памяти и реакция на внешние события.

Способ реализации заметно влияет на образовательный эффект. Платформа может быть описана на языках аппаратной логики (Verilog, VHDL, SystemVerilog) и предназначаться для синтеза на ПЛИС, что даёт потенциальную возможность физического воплощения. Альтернативой служит программный эмулятор на языках высокого уровня (Python, C, Java), представляющий собой чисто абстрактную модель без прямой связи с аппаратурой. Первый подход обеспечивает наглядность «железной» стороны архитектуры, второй — быстроту запуска и простоту модификации.

Наконец, интерфейс и интерактивность определяют удобство работы. Наличие текстовой консоли (ассемблерная консоль, командная строка) и/или графического интерфейса с визуализацией регистров, памяти, шин, поддержкой пошагового выполнения и отладки существенно усиливает обучающий эффект, делая поведение модели управляемым и понятным.

Существующие учебные платформы можно расположить на спектре от простейших абстракций до синтезируемых процессорных ядер, способных исполнять полноценную операционную систему.

Модель «Маленького человечка» (Little Man Computer, LMC), предложенная Стюартом Мэдником в 1965 г. \cite{LMC}, представляет аналогию почтового отделения: в закрытой комнате «человечек» выполняет инструкции, манипулируя данными в ячейках памяти. Память содержит 100 ячеек (адреса 0–99), каждая из которых хранит трёхзначные десятичные числа. Ввод и вывод осуществляются через воображаемые лотки INBOX и OUTBOX, арифметические операции ограничены сложением и вычитанием. LMC отлично подходит для начальной демонстрации принципа хранимой программы, однако не даёт представления о конвейере, иерархии памяти или взаимодействии с периферией.

Архитектура SIC/XE (Simplified Instructional Computer Extra Equipment) \cite{SICXE} расширяет базовую модель SIC и применяется в ряде учебников по системному программированию. Объём памяти увеличен до 1 Мбайт, организованной 8‑битными байтами; машинное слово состоит из трёх последовательных байтов (24 бита). Девять регистров, включая базовый и регистры с плавающей запятой, а также четыре формата инструкций и несколько режимов адресации (прямая, индексированная, базовая) приближают модель к реальным процессорам. Набор инструкций дополнен операциями с плавающей запятой и системными вызовами для обработки прерываний. Однако реализация существует лишь в форме программного симулятора, что не позволяет продемонстрировать связь с физической аппаратурой.

Модель «Е14» \cite{E14} ориентирована на изучение многоядерной архитектуры и параллельного программирования. В ней реализована несимметричная конфигурация, где один из процессоров выполняет функции ведущего. Система команд совместима с предшествующей моделью «Е97», что упрощает переход от однопроцессорных концепций к многопроцессорным. Межпроцессорный обмен использует порты, общую шину и механизмы прямого доступа к памяти. «Е14» работает под управлением Windows как прикладная программа и не рассчитана на запуск операционной системы, что ограничивает её применение при изучении системного ПО.

Отечественная учебная модель ЭВМ \cite{ourcomputer} делает акцент на работе кэш-памяти. В её состав входят процессор, оперативная память и сверхоперативная память, объединяющая регистры общего назначения и кэш, а также устройства ввода/вывода. Модель наглядно демонстрирует ускорение доступа к данным, но не затрагивает верхние уровни организации вычислительного процесса.

Проект Nand2Tetris \cite{Nand2Tetris} реализует комплексный подход, проводя обучающегося через всю иерархию вычислительной системы: от логического вентиля NAND до операционной системы и компилятора. Аппаратная часть базируется на 16‑битной архитектуре Hack с 17 инструкциями и поддержкой экранного вывода и клавиатуры. Программная часть включает ассемблер, виртуальную машину со стековой организацией и компилятор объектно-ориентированного языка Jack. Несмотря на глубину проработки, полное прохождение курса требует более недели практических занятий, что трудно вписать в типовой семестровый график. Кроме того, использование собственной архитектуры и языков программирования изолирует получаемые навыки от индустриальных стандартов. Интерактивная игра Nand Game, решающая сходную задачу в упрощённой форме (построение вентилей, ALU, регистров, памяти на D‑триггерах и процессора начиная с NAND-вентиля), даёт лишь поверхностное представление и недостаточна для университетского курса.

На другом полюсе находятся синтезируемые ядра промышленного уровня. Ядро RISC‑V CPU Core (RV32IM) UltraEmbedded представляет собой 32‑битный процессор на Verilog, поддерживающий набор инструкций RV32IM, пользовательский, супервизорный и машинный режимы, систему управления памятью (MMU), эмуляцию атомарных операций и интерфейсы AXI \cite{UltraEmbedded}. Оно позволяет запускать ОС Linux и проверено с помощью средств co‑simulation и Google RISCV‑DV. Однако полнота и сложность этого ядра делают его скорее объектом изучения готовой архитектуры, нежели платформой для поэтапного освоения принципов проектирования.

Таким образом, ни одна из рассмотренных учебных платформ не объединяет открытый процессор со стандартной RISC‑архитектурой, легковесную UNIX‑подобную ОС и компактный компилятор, обеспечивая при этом демонстрацию цикла разработки и выполнения программ. Этот пробел обосновывает необходимость создания нового программного комплекса.

\subsection{Выводы и требования к учебному стенду}

Обзор открытых процессорных архитектур показал, что RISC‑V RV32I является оптимальной основой для учебного процессора благодаря свободной лицензии, модульности и развитой экосистеме. Среди легковесных UNIX‑подобных ОС выделяется Xv6: она переписана на ANSI C, хорошо документирована, содержит около 9000 строк кода и поддерживает архитектуру RISC‑V, что делает её подходящим объектом для изучения механизмов управления процессами, памятью и вводом‑выводом. В качестве компилятора языка Си обоснован выбор TinyCC, сочетающего умеренный объём исходного кода с гибкостью адаптации бэкенда. Анализ учебных компьютерных платформ выявил отсутствие решений, интегрирующих все три компонента в единый стенд.

На основании проведённого анализа формулируются следующие требования к проектируемому комплексу:

\begin{enumerate}
  \item Процессорное ядро должно реализовывать спецификацию RISC‑V RV32I (базовый целочисленный набор инструкций), быть синтезируемым на ПЛИС и обеспечивать генерацию сигналов трассировки.
  \item Операционная система — Xv6, собранная для целевой архитектуры и способная выполняться на виртуальном процессоре написанном на verilog.
  \item Компилятор — TinyCC интегрированный в сборочную цепочку.
  \item Комплекс обязан реализовывать маршруты: компиляция программы и ее запуск на виртуальном процессоре; запуск xv6 на виртуальном процессоре и компиляция программы внутри с помощью TinyCC.
  \item Для наглядности необходимо предоставить средства трассировки (VCD‑файлы) и отладки состояния процессора, памяти и периферийных устройств.
  \item Воспроизводимость платформы должна обеспечиваться унифицированным сборочным окружением (Docker‑образ) и едиными репозиториями всех компонентов.
  \item Комплекс сопровождается учебно‑методическими материалами, описывающими архитектуру, процедуру сборки и сценарии лабораторных работ.
\end{enumerate}

Детальная архитектура стенда, отвечающего перечисленным требованиям, рассматривается в следующем разделе.
